\documentclass[12pt]{article}
\usepackage[utf8]{inputenc}
\usepackage[T1]{fontenc}
\usepackage{hyperref}
\title{Carútana Language Resources for K-12 Curriculum Development}
\author{Compiled by Claude Sonnet 5.5 (claude-sonnet-5-5) for Andr\'as Kornai}
\date{}
\begin{document}
\maketitle

\section{Introduction}
Carútana (ISO 639-3 \texttt{cru}) is treated in the sources found as a variety of Baniwa (Karu), a Northern Arawakan language of the upper Rio Negro region of Brazil, Colombia and Venezuela. Sources disagree on whether Carútana is a dialect or a separate language, and one source states that it is extinct. The identification supplied with the request (family ``Bookkeeping'') is a Glottolog placeholder, not a genealogical classification. The searches found no Carútana-specific educational material.

\section{Primary Educational Resources}
No primers, textbooks or school programmes specific to Carútana were found in the search results.

\section{Linguistic and Reference Materials}
\begin{itemize}
\item Wikipedia, Baniwa of I\c{c}ana (the Baniwa variety closest to the Carútana context): \url{https://en.wikipedia.org/wiki/Baniwa_of_I%C3%A7ana}
\item Wikipedia, Karu language (discusses the classification of the Baniwa/Karu varieties): \url{https://en.wikipedia.org/wiki/Karu_language}
\item Endangered Languages Project record: \url{https://endangeredlanguages.com/node/88596}
\end{itemize}

\section{Religious and Community Resources}
\begin{itemize}
\item Joshua Project, Carutana in Brazil (people-group profile): \url{https://joshuaproject.net/index.php/people_groups/19321}
\end{itemize}

\section{Digital and Audio Resources}
None specific to Carútana were found.

\section{Curriculum Development Context}
Little exists under the name Carútana. Any curriculum work would probably have to start from Baniwa materials, after confirming with the community which variety is meant and whether it is still spoken.

\end{document}